\documentclass{internal-report}
\usepackage[UTF8]{ctex}
% STATUS: done
% LAST_MODIFIED: 2026-08-21
%
% ============================================================================
% S3 内部报告 · 内容段（iter-02-sub3-cross-disease）
% 阶段 2.3 内容段：门禁 2 后取数填数字（S3-results.pkl 落盘）。
% 所有数字均从 outputs/data/S3-results.pkl 只读提取（过滤后特征数/CLR 常数
% 从 S3-preprocessed.pkl 只读验证），与 handoff-S3-model-agent.md 一致。
% 正文写人话不标 pkl 字段路径；溯源进 writing-material-S3.tex。
% 图表占位符保留，由队友/代码子代理出图后替换。
% ============================================================================

% ----------------------------------------------------------------------------
% 表述库映射（骨架指南，供内容段取用；正式写作时逐句落位，见
% knowledge/expression-library.md 对应模板 ID）
%   §1 问题与目标      → 表述库「引入问题」：数据现象切入式 / 必要性论证式
%   §2 数据与 LODO 评估框架 → 表述库「摘要-数据处理说明」「模型引入段-过渡到模型」：
%                         发现式预处理（近全零过滤→CLR→标准化）+ 评估框架自足定义
%   §3 四策略构建与对比 → 表述库「模型引入段-方法比较式 / 特点+适用式」：
%                         逐策略「动机句→构建方式→测试集信息隔离边界」
%   §4 最优策略交付     → 表述库「结果与分析段-对比择优式 / 直接汇报最优解式」
%   §6 性能探讨         → 表述库「结果与分析段-量化结论+规律解释式」、
%                         「模型评价段-缺点写法（保留真实局限）」
%   §7 结论             → 表述库「模型评价段-改进与推广式」
% 数字标准：正文写人话不标 pkl 字段，溯源进写作素材独立文件
%           writing-material-S3.tex（内容段生成）。
% ----------------------------------------------------------------------------

\title{S3 内部报告：跨疾病预测模型\\（子问题 3，leave-one-disease-out 泛化评估）}
\author{报告对话（S3）}
\date{2026-08-21}

\begin{document}

\maketitle

\begin{abstract}
\textbf{内容摘要：}本报告回答题面 S3「建立一个跨疾病预测模型，并探讨模型在跨疾病预测上的性能」。核心叙事：本问解决「模型到新疾病上还灵不灵」→ 建立 LODO（留一疾病）诚实评估框架 → 正式构建并对比四种跨疾病建模策略 → 选出最优策略作为交付模型 → 若全部策略性能不达标则触发后续模型探索流程重新选模型 → 从衰减归因 / 深度迁移 / 阈值漂移三视角探讨跨疾病性能 → 结论。\textbf{关键结果}：四策略 3 组合 AUC 均值全部低于 0.60（最优策略 A 直接迁移均值 0.5603），触发后续模型探索流程；候选模型序列 R1--R4（候选模型族，见 §5）中密度比重加权域适应（R3）达到最优可达均值 0.6068，仍低于可用性参考线 0.65，故本问以「尝试后的完整验证过程」交付，未交付可用模型。衰减归因显示主导来源为疾病特异信号（IBD 衰减最大 $-0.271$）与患病定义差异（Obesity）；深度迁移分析显示共享物种方向一致性接近随机（51.2\%，符号检验 $p=0.536$）；C3 阈值漂移量化显示训练/测试患病基线差 $\Delta=+0.332$，Youden 阈值落在测试分布 96.0\% 分位，灵敏度仅 0.024。
\end{abstract}

% 文首符号表（自足性规则；正文首次出现处均有内嵌定义）
\begin{table}[ht]
\centering
\caption{符号表（本文自足性速查；每个符号正文首次出现处均有内嵌定义）}
\begin{tabular}{c p{0.50\textwidth} c}
\toprule
\textbf{符号} & \textbf{定义} & \textbf{首次出现} \\
\midrule
LODO & 留一疾病（leave-one-disease-out）：每次留一种疾病作测试集，其余作训练集，共 3 组合 & §2 \\
C1 / C2 / C3 & 三个 LODO 组合（测试疾病分别为 CRC / IBD / Obesity） & §2 \\
AUC & 曲线下面积（主指标，阈值无关） & §3 \\
$J$ & Youden J 统计量（灵敏度$+$特异度$-1$），$\tau^*$ 为其最优阈值 & §3 \\
$P_{\mathrm{cal}}$ & Platt 校准概率，参数 $A,B$（仅训练集估计） & §3.4 \\
silhouette & 轮廓系数（$[-1,1]$，衡量聚类分离度） & §6.1 \\
\bottomrule
\end{tabular}
\end{table}

\begin{table}[ht]
\centering
\caption{跨问 / 裁定代号表}
\begin{tabular}{c p{0.50\textwidth} c}
\toprule
\textbf{代号} & \textbf{指代} & \textbf{首次出现} \\
\midrule
S3 门禁 1 重构 & 主分析裁定（2026-08-21）：四策略对比升级为正式建模 + 后续模型探索流程 & §1 \\
S1 标准 & 策略 A 沿用的 S1 已验证模型标准（L2+CLR） & §3.1 \\
S2 生物标志物 & 子问题 2 的生物标志物识别结论，本问在 §6/§7 交叉核对 & §6.2 \\
三分法归因 & 批次效应 / 疾病特异信号 / 患病定义差异（互斥归因来源） & §6.1 \\
\bottomrule
\end{tabular}
\end{table}

% ============================================================================

\section{问题与目标}

本节约一句话：本问解决「模型到一个没见过的疾病上还灵不灵」，正面回答是「建立跨疾病预测模型」——正式构建并对比四种策略、选出最优者交付，性能不达标则走后续模型探索流程，直到交付可用模型或候选模型均已尝试后以完整验证过程报告。

\subsection{本问解决什么}
题面 S3 要求建立跨疾病预测模型并探讨其性能。这里的核心是「跨疾病」：模型在若干疾病上学到的判别能力，能否迁移到训练阶段从未见过的另一种疾病。本文不把「负结果」当作默认终点，而是按主分析裁定（\texttt{S3 门禁 1 重构}）先尽力交付模型，失败验证过程是均已尝试后的最后结论。

\subsection{怎么解决（故事线总览）}
回答按以下链路推进：
\begin{enumerate}
  \item 建立 LODO（留一疾病）评估框架，作为「新疾病不可见」的诚实评估基准（§2）；
  \item 正式构建并对比四种跨疾病建模策略（A 直接迁移 / B 共享标志物 / C 分类学聚合 / D 输出概率修正），选出最优者（§3）；
  \item 交付最优策略，报告其性能水平与适用条件（§4）；
  \item 若四策略性能均不达标，触发后续模型探索流程重新选模型，记录过程与交付（§5，条件节）；
  \item 从衰减归因 / 深度迁移 / 阈值漂移三视角探讨跨疾病性能为什么好或为什么差（§6）；
  \item 结论收束（§7）。
\end{enumerate}

\subsection{主/次目标声明}
本问为单一评估目标（度量跨疾病预测性能），AUC 为唯一主指标，辅指标是 AUC 的分解视图，不构成独立的次目标，因此不写主次目标取舍结论。

% ============================================================================

\section{数据与 LODO 评估框架}

本节约一句话：交代数据来源与预处理标准，并正式定义 LODO 三组合评估框架——测试疾病在训练阶段完全不可见。

\subsection{数据}
使用 0.3 清洗后的物种级相对丰度数据（定和成分数据）。预处理与 S1 标准一致、严格测试集信息隔离（S1：子问题 1，单疾病患病预测模型，本问沿用其已验证的 L2+CLR 预处理标准；S2：子问题 2，生物标志物识别，本问在 §6/§7 与其结论交叉核对）：
\begin{enumerate}
  \item 近全零过滤（零值占比超过阈值剔除，三病并集统一标准，与 S1/S2 一致）——过滤后特征数\ 264（由原始 1331 个物种特征过滤得到，样本共 484 例）；
  \item CLR 变换（成分数据强制，零值乘法替换）——替换常数\ $6.5\times10^{-6}$；
  \item StandardScaler（均值/方差\textbf{仅训练集估计}，测试集信息隔离）。
\end{enumerate}

\subsection{LODO 评估框架}
本节说明为什么用 LODO：它是「跨疾病」最诚实的评估方式——测试疾病在训练阶段完全不可见（标签与特征均不参与训练），模型学到的是真正可泛化的判别。

三个组合（训练集为其余两种疾病）：
\begin{enumerate}
  \item \textbf{C1}：测试 \texttt{CRC}，训练 \texttt{IBD + Obesity}，测试集正类占比\ 39.7\%；
  \item \textbf{C2}：测试 \texttt{IBD}，训练 \texttt{CRC + Obesity}，测试集正类占比\ 22.7\%（少数类）；
  \item \textbf{C3}：测试 \texttt{Obesity}，训练 \texttt{CRC + IBD}，测试集正类占比\ 64.8\%（正类占多数）。
\end{enumerate}

% ============================================================================

\section{四策略构建与对比}

本节约一句话：正式构建四种跨疾病建模策略（直接迁移 / 共享标志物 / 分类学聚合 / 输出概率修正），在同一 LODO 评估框架下对比 AUC，为选出最优交付模型提供证据。

每个策略的写法统一为「动机句（为什么看这个）→ 构建方式 → 测试集信息隔离边界 → 预期观测」。

\subsection{策略 A：直接迁移（基线）}
\textbf{动机}：先看最直接的跨疾病做法——把在训练疾病上训练好的模型原样搬到测试疾病，作为全部策略的基线。
\textbf{构建}：正则化 Logistic 回归（L2 惩罚，$C=1.0$（$C$：L2 正则化强度超参数），class\_weight balanced）+ CLR 前置 + 近全零过滤物种级特征（264 维），沿用 S1 标准。
\textbf{评估}：主指标 AUC（阈值无关，最诚实）；辅指标为训练集 Youden J 最优阈值迁移下的 ACC（准确率）/ 灵敏度 / 特异度 / F1（标注测试集类别不平衡）\cite{youden1950}。
\textbf{测试集信息隔离}：\textbf{禁止在测试集上重定阈值}（会用到新疾病标签）。

\subsection{策略 B：跨疾病通用模型（共享标志物特征子集）}
\textbf{动机}：若疾病特异信号在拖后腿，则只保留三疾病共享的物种（疾病无关特征）重训「通用患病模型」，看能否恢复性能。
\textbf{构建}：取三数据集共享物种（基于过滤后特征集按特征名取交集，数量\ 252）作特征子集，在训练疾病上重训策略 A 同款模型，重跑 LODO。
\textbf{测试集信息隔离}：特征交集仅使用\textbf{特征存在性}（哪些物种在三数据集都出现），\textbf{绝不用测试集标签}（仅使用测试集特征而不使用标签显式声明——指特征选择/权重估计等环节不使用测试集标签，仅用特征存在性等无监督信息）。

\subsection{策略 C：分类学聚合模型（属级/门级 + CLR）}
\textbf{动机}：回应题面注 2「可利用微生物的分类学信息」——把物种级特征聚合到属级/门级，看分类学层级能否带来增益。
\textbf{构建}：物种级特征按属（\texttt{g\_\_}）/ 门（\texttt{p\_\_}）层级聚合（同属/同门丰度求和，聚合后维度大幅下降：属级 106 维、门级 11 维），CLR 后重训同款模型，重跑 LODO。
\textbf{预期观测}：聚合是否提升跨疾病性能为\emph{实测观测}，如实报告（分类学信息已利用，是否有增益不预判）。

\subsection{策略 D：输出概率修正模型（最优策略 + 概率校准 + 阈值迁移）}
\textbf{动机}：AUC 与阈值无关、不受校准影响，但可直接应用的水平（ACC/F1/灵敏度）会因训练与测试患病概率基线不同而失准。本策略在最优策略上叠加 Platt 校准，修复可直接应用的指标。
\textbf{构建}：在策略 A/B/C 中 AUC 最优者（实测为策略 A）上叠加 Platt 缩放\cite{platt1999}，把原始分数映射到校准概率，再迁移 Youden 阈值。
\textbf{测试集信息隔离}：校准参数 $A,B$（$P_{\mathrm{cal}}(y{=}1\mid\mathbf{x}) = 1/(1+\exp(A\,f(\mathbf{x})+B))$）\textbf{仅训练集估计}；实现校验 $A<0$（等价 $w>0$，防排序反转）。\textbf{不改 AUC}（单调变换不改变排序），只修复可直接应用的指标。

\subsection{策略对比（核心交付）}

以下为四策略 3 组合 AUC 对比表。判定（各策略是否达可用性参考线）见 §5——四策略均值全部低于 0.60，触发后续模型探索流程。

\begin{table}[ht]
\centering
\caption{四策略 3 组合 AUC 对比（核心交付；C 分属级/门级两档，D 以策略 A 为基）}
\begin{tabular}{l c c c c}
\toprule
\textbf{策略} & \textbf{构建方式} & \textbf{C1 AUC} & \textbf{C2 AUC} & \textbf{C3 AUC} \\
\midrule
A 直接迁移 & L2+CLR 物种级 264 & 0.5674 & 0.5882 & 0.5253 \\
B 共享标志物 & 共享特征子集 252 重训 & 0.5417 & 0.6080 & 0.5218 \\
C 属级聚合 & 属级 106 维 + CLR & 0.3616 & 0.4861 & 0.5440 \\
C 门级聚合 & 门级 11 维 + CLR & 0.4141 & 0.5261 & 0.5999 \\
D 输出概率修正 & 策略 A + Platt 校准 & 0.5674 & 0.5882 & 0.5253 \\
\bottomrule
\end{tabular}
\label{tab:strategy-compare}
\end{table}

四策略 3 组合 AUC 均值：A 0.5603、B 0.5572、C 属级 0.4639、C 门级 0.5134、D 0.5603（与基策略 A 相同，符合 Platt 单调变换不改排序）。\textbf{全部低于 0.60}，触发后续模型探索流程（§5）。

% 图表规格：
%   数据源: S3-results.pkl strategy_compare.{A_direct,B_shared,C_genus,C_phylum,D_calibrated}.*.auc + mean_auc
%   图名: 四策略跨疾病 AUC 对比（表）
%   论文位置: 结果与讨论（策略对比节）

% ============================================================================

\section{最优策略交付}

本节约一句话：从四策略中选定交付模型，报告其构建方式、跨疾病性能水平与适用条件；若后续流程被触发，则此节交付探索所得模型（见 §5）。

\subsection{交付模型选择}
选定依据（策略对比表 \ref{tab:strategy-compare} + 可用判定标准）——四策略均值全部低于可用性参考线，故交付模型由候选模型序列决定（§5），最优可达策略为 \textbf{R3 密度比重加权域适应}（3 组合 AUC 均值 0.6068）。
选定理由与备选取舍：四策略中策略 A 直接迁移均值最高（0.5603），但低于可用性参考线 0.65；候选模型序列中 R3 加权域适应达到最优可达均值 0.6068，仍低于可用性参考线，故本问以「最优可达 + 完整验证过程」交付，不宣称可用模型。

\subsection{性能水平与适用条件}
- 3 组合 AUC 均值：\ 0.6068（R3 加权域适应）；各组合取值：C1 0.5945、C2 0.6489、C3 0.5771。
- 阈值迁移辅指标（R3，训练集 Youden 阈值迁移）：C1 ACC 0.6446 / 灵敏度 0.2292 / 特异度 0.9178 / F1 0.3385；C2 ACC 0.5727 / 灵敏度 0.6800 / 特异度 0.5412 / F1 0.4198；C3 ACC 0.4150 / 灵敏度 0.1585 / 特异度 0.8876 / F1 0.2600。
- 适用条件与实际应用建议：\ 本问未交付可用模型（最优可达 0.6068 < 可用性参考线 0.65）。R3 加权域适应在 C2（IBD）上表现相对最好（0.6489），在 C3（Obesity）上仍受阈值漂移拖累（灵敏度 0.1585）。应用建议：若必须实际应用，优先用于患病基线接近训练分布的疾病，并需在目标疾病上重定阈值。
- 真实方法局限（保留精确边界）：\ 域适应权重基于域分类器估计（Logistic 区分 train/test），权重裁剪上界 10；R3 提升相对策略 A 为 $+0.0465$，未达「相对 A 提升 $\ge 0.10$」的可用性参考线。

% ============================================================================

\section{后续模型探索记录（触发时）}

本节约一句话：四策略 3 组合 AUC 均值均低于可用性参考线，触发后续模型探索流程重新选模型；本节记录触发证据与逐级尝试结果，直到交付可用模型或候选模型均已尝试。

\subsection{触发判定}
触发条件（四策略均值均\ $<$\ 0.60）：实测四策略均值全部低于 0.60，\textbf{触发}。可用性参考线定义为「候选模型均值 $\ge 0.65$ 或相对策略 A 提升 $\ge 0.10$」。

\subsection{逐级尝试}
候选模型族按序尝试，每级同 LODO 评估框架评估，达可用性参考线即交付：
\begin{enumerate}
  \item \textbf{R1 树模型族}（Random Forest 500 树，物种级过滤后特征）——均值 0.5092，未达可用性参考线；
  \item \textbf{R2 样本合并通用模型}（2 训练疾病样本合并训练「患病 vs 健康」通用分类器）——均值 0.5603（与策略 A 数学恒等，因 LODO 评估框架本身已合并 2 训练疾病），未达可用性参考线；
  \item \textbf{R3 密度比重加权域适应}（importance weighting，域分类器法，仅使用测试集特征而不使用标签显式声明）——均值 0.6068，未达可用性参考线（$<0.65$，相对 A 提升 $+0.0465<0.10$）；
  \item \textbf{R4 对抗式域适应}（DANN，即 Domain-Adversarial Neural Network 域对抗神经网络，梯度反转，最后手段，严格验证集测试集信息隔离）——均值 0.5947，未达可用性参考线。
\end{enumerate}

\subsection{完整验证过程}
R1--R4 全部未达可用性参考线，报告最优可达性能 + 失败主因 + 实际应用建议（完整验证过程）：
- 最优可达性能：\ R3 加权域适应 3 组合 AUC 均值 0.6068，低于可用性参考线 0.65；
- 失败主因：\ 衰减归因显示主导来源为疾病特异信号（IBD 衰减最大）与患病定义差异（Obesity），共享物种方向一致性接近随机（§6）；
- 实际应用建议：\ 本问不交付可用模型，以验证过程说明「跨疾病泛化在当前数据与框架下受限」，并给出改进方向（§7）。

% ============================================================================

\section{性能探讨}

本节约一句话：从衰减归因（三分法）、深度迁移（什么在迁移）、阈值漂移（C3 诊断）三视角解释跨疾病性能为什么是这个水平，为最优策略的可用边界与后续改进提供依据。

\subsection{衰减归因：三分法}
\textbf{动机}：把每个低 AUC 组合的失败归因到三个\textbf{互斥}来源，每个来源配可计算的定量证据，避免笼统说「泛化差」。
三分法来源与操作化定义：
\begin{enumerate}
  \item \textbf{批次效应}：训练/测试来自不同研究/平台导致的分布偏移——定量证据 silhouette 系数（数据集标签）\ $=0.070$（A 类验证值，近 0，不主导）；
  \item \textbf{疾病特异信号}：模型学的是训练疾病特有标志物——定量证据域内 AUC 减跨疾病 AUC 的衰减量（见下表）；
  \item \textbf{患病定义差异}：新疾病患病概率基线整体偏移、决策边界失准——定量证据训练/测试患病基线差（§6.3）。
\end{enumerate}

衰减归因表（各疾病域内/跨疾病 AUC、衰减量、主导归因；域内 AUC 为 264 特征 5 折 CV（5-fold cross-validation，五折交叉验证）重算，与跨疾病同标准）。\textbf{标准声明}：本问预处理含 StandardScaler（均值/方差仅训练集估计），与 S1/S2 的纯 CLR 标准不同，故域内 AUC 与 S1 有差异（$\Delta$ $-0.010$/$-0.028$/$+0.014$，CRC/IBD/Obesity）：
\begin{table}[ht]
\centering
\caption{衰减归因表（三分法；衰减量 = 跨疾病 AUC $-$ 域内 AUC）}
\begin{tabular}{l c c c c}
\toprule
\textbf{疾病} & \textbf{域内 AUC} & \textbf{跨疾病 AUC} & \textbf{衰减量} & \textbf{主导归因} \\
\midrule
CRC & 0.7811 & 0.5674 & $-0.2138$ & 疾病特异信号 \\
IBD & 0.8588 & 0.5882 & $-0.2706$ & 疾病特异信号（最强） \\
Obesity & 0.6638 & 0.5253 & $-0.1384$ & 患病定义差异 \\
\bottomrule
\end{tabular}
\label{tab:decay-attribution}
\end{table}

% 图表规格：
%   数据源: S3-results.pkl decay_attribution.{CRC,IBD,Obesity}.{domain_auc,cross_auc,decay,dominant_cause}
%   图名: 三分法衰减归因表
%   论文位置: 结果与讨论（衰减归因节）

\subsection{深度迁移分析：什么在迁移、什么不迁移}
\textbf{动机}：回答「哪些共享标志物的丰度变化方向跨疾病一致（可迁移信号）、哪些翻转（疾病特异信号）」，把共享物种分成两类并量化占比，识别真正可迁移的生物学信号。
- 方向一致（可迁移）数量与占比：\ 387 个，占 51.2\%（共 756 个有效共享物种）；
- 方向翻转（疾病特异）数量与占比：\ 369 个，占 48.8\%；
- 符号检验：\ $p=0.5364$（不显著）——共享物种的「患病↑/↓」方向跨疾病接近随机（$\approx$50/50），即共享物种的\textbf{存在性}不承载可迁移信号，信号藏在\textbf{疾病特异方向}里；
- 可对接 S2 生物标志物结论的清单：\ 见写作素材（共享物种清单 252 个，方向一致性逐物种标注，供 S2 生物标志物交叉核对）。

\begin{figure}[ht]
\centering
% 图占位：迁移方向一致性图（数据已就绪，图待队友/代码子代理出图后替换）
% 图表规格：
%   数据源: S3-results.pkl migration_analysis.{direction_consistent_count,direction_flipped_count,shared_species_list}
%   图名: 共享标志物跨疾病迁移方向一致性（方向一致 vs 方向翻转）
%   论文位置: 结果与讨论（深度迁移分析节）
\caption{共享标志物跨疾病方向一致性示意——数据已就绪（387 一致 / 369 翻转），图待出图后替换}
\label{fig:migration-direction}
\end{figure}

\subsection{阈值漂移量化（C3）}
\textbf{动机}：把「AUC 接近随机而可直接应用指标崩溃」的矛盾展开成完整诊断——训练集与测试集患病概率基线差多少、Youden 阈值迁移后决策边界落在测试集分布何处。
- 训练集 vs 测试集患病概率基线差：\ 训练 0.316 vs 测试 0.648，基线差 $\Delta=+0.332$；
- Youden 阈值迁移后决策边界在测试集分布的位置：\ $\tau^*=0.9205$ 落在测试分数分布的 96.0\% 分位（96.0\% 测试样本被判健康）；
- 完整诊断（AUC 与可直接应用指标为何背离）：\ 测试分布整体左移，训练集拟合的 Youden 阈值迁移后决策边界失准，灵敏度仅 0.024（几乎全判健康）——患病定义差异主导，AUC 接近随机（0.525）而可直接应用指标崩溃。Platt 校准为单调变换，不改变决策边界，故无法修复 C3 可直接应用指标（即使改用 0.5 概率阈值，灵敏度仅 0.165）。

\begin{figure}[ht]
\centering
% 图占位：阈值漂移诊断图（数据已就绪，图待队友/代码子代理出图后替换）
% 图表规格：
%   数据源: S3-results.pkl threshold_drift.{train_baseline,test_baseline,boundary_position,diagnosis}
%   图名: C3 阈值漂移诊断（训练/测试患病概率基线 + 决策边界位置）
%   论文位置: 结果与讨论（阈值漂移节）
\caption{C3 阈值漂移诊断示意图——数据已就绪（基线差 $\Delta=+0.332$，阈值落 96.0\% 分位），图待出图后替换}
\label{fig:threshold-drift}
\end{figure}

% ============================================================================

\section{结论}

本节约一句话：收束本问对「建立跨疾病预测模型」的正面回答——交付模型 + 性能水平 + 适用条件 + 局限与改进方向。

\begin{itemize}
  \item 交付模型与跨疾病性能：\ 本问未交付可用模型。四策略均值全部低于 0.60（最优策略 A 直接迁移 0.5603），触发后续模型探索流程；候选模型序列 R1--R4 中 R3 密度比重加权域适应达到最优可达均值 0.6068，仍低于可用性参考线 0.65，以「最优可达 + 完整验证过程」交付。
  \item 性能为何是这个水平（衰减/迁移/漂移主因）：\ 衰减归因显示主导来源为疾病特异信号（IBD 衰减最大 $-0.271$）与患病定义差异（Obesity）；深度迁移分析显示共享物种方向一致性接近随机（51.2\%，$p=0.536$），共享物种存在性不承载可迁移信号；C3 阈值漂移量化显示训练/测试患病基线差 $\Delta=+0.332$，Youden 阈值落测试分布 96.0\% 分位，灵敏度仅 0.024。
  \item 适用条件与真实局限：\ 若必须实际应用，优先用于患病基线接近训练分布的疾病，并需在目标疾病上重定阈值；域适应权重基于域分类器估计（裁剪上界 10），R3 相对 A 提升 $+0.0465$ 未达可用性参考线。
  \item 改进与推广方向：\ (1) 引入更多疾病/更大样本以缓解疾病特异信号主导；(2) 针对患病定义差异做目标疾病阈值重定或校准；(3) 探索疾病特异方向信号（而非共享存在性）作为可迁移特征；(4) 与 S2 生物标志物结论交叉核对，识别真正可迁移的生物学信号。
\end{itemize}

% 写作素材见 writing-material-S3.tex（独立文件，供队友阶段 4 写作）
% READER-AUDIT: 见 reader-audit-sub3.md

\begin{thebibliography}{9}
\bibitem{platt1999}
Platt, J. C. Probabilistic outputs for support vector machines and comparisons to regularized likelihood methods. In \emph{Advances in Large Margin Classifiers}, 61--74, 1999.
\bibitem{youden1950}
Youden, W. J. Index for rating diagnostic tests. \emph{Cancer}, 3(1):32--35, 1950.
\end{thebibliography}

\end{document}
